\begin{abstract}
Temporal validity and sequential detection operate on different statistical
clocks under drift. For finite-memory distribution tracking, a fixed operating
window can cross the temporal-validity horizon before standard changepoint
detectors accumulate enough evidence to alarm under false-alarm control. The
object of this paper is the delay
\[
  \Delta_{\mathrm{val-det}} \coloneqq \tau_{\mathrm{detect}} - \tau_{\mathrm{valid}}(n),
\]
where $\tau_{\mathrm{valid}}(n)$ is the expiry time of an operating memory level
$n$ and $\tau_{\mathrm{detect}}$ is the detector stopping time on the same
stream. The target theory is a detector-class sign law, and ultimately a lower
bound, showing that positive pre-detection invalidity is structural for H\"older
drift families under false-alarm constraints. The benchmark proof route starts
from a one-dimensional Gaussian location model with H\"older mean path and uses
the temporal-validity horizon from finite-memory tracking to define the expiry
clock. Compact false-alarm-calibrated experiments already support the target:
observation-based ADWIN, Page--Hinkley, and CUSUM retain positive
validity-detection gaps across tested roughness exponents $H\in\{0.5,0.75,1.0\}$,
while residual-input variants and KSWIN often trade larger delays for missed
alarms.
\end{abstract}
